\honest{The Sheer Folly of Callow Youth}{Eliezer Yudkowsky}{2008}

I have hinted that I once gave a mysterious answer to a mysterious question. It was not only about consciousness: ``the more embarrassing error was that I took a mysterious view of morality.'' \nb{In ``My Wild and Reckless Youth'' Yudkowsky had declined to describe the answer, and gave only its prediction: that neurons exploit quantum gravity, after Roger Penrose.}

In 1997 I set out to prove that building a superintelligence was right. Does life have meaning? ``I don't know,'' I said, and then ``But.'' If life has no meaning, every outcome has utility zero, so that case cancels out and I can act as if life is meaningful. Humans cannot find the meaning, so a superintelligence must. I do not trust my memory of this, and I will not reread my old writings. \nb{On every point checked, the reconstruction matches Yudkowsky's ``Frequently Asked Questions about the Meaning of Life'' (1999; the copy checked was last updated in 2001), down to its conclusion that the Singularity is ``the interim meaning of life.''}

What went wrong? I treated utility as a kind of stuff, so ``life is meaningless'' meant zero; but a world where everything is equally meaningful would drop out just the same. \nb{A recent book on moral uncertainty, by MacAskill, Bykvist and Ord, makes the same point: a theory on which all options are equally good ``never affects what it's appropriate to do,'' whatever number it gives them.} Worse, I took ``should'' to be whatever compels any mind, and assumed that this is what people mean by ``right.'' \nb{The book argued against that equation in ``No Universally Compelling Arguments.''} And having ``proved'' the meaningless case irrelevant, I saw no need to define ``intelligence'' or ``meaning.''

My lesson: ``No matter how clever the justification for relaxing your standards, or evading some requirement of rigor, it will blow your foot off just the same.'' \nb{The rule is stated for every case. The evidence offered is this one.} I was skilled at refuting others and never turned that skill on myself.

I wanted speed, because of the Singularity and the danger of nanoweapons. ``No, you don't use the best concepts you can use at the time.'' Nature accepts no excuses: ``Either you match it, or you fail.'' Then I narrow the rule: you may think with vague concepts, but you do not build a superintelligence on them. \nb{The premise, that a best guess fails as surely as no attempt, is asserted here. The book argues it in ``Value is Fragile,'' written four months later.}

It gets worse. To the objection that no ``ought'' follows from pure logic, I replied that three ``hard problems,'' consciousness, existence and morality, were linked, too hard for humans, and probably needed new physics. \nb{The objection is Hume's. The physics of consciousness is Roger Penrose's.}

So: ``You cannot manipulate confusion.'' The alchemists could not know why lead would not turn into gold, ``So they poisoned themselves and died.'' \nb{The deaths from alchemy on record that the notes found came from elixirs meant to prolong life, not from attempts to make gold.}
